\subsection*{一步一步建立直观}

\paragraph{1. 寻找增长速度等于当前存量的刻度}

选择指数函数。把纵轴读作某个不断增长的存量，切线斜率就是瞬时增长速度。自然指数的特殊性在于：把时间单位选定后，任一点的增长速度恰好等于该点的存量。这是一条局部演化规则，而不是从一个小数近似出发的定义故事。



\paragraph{2. 对照不同函数，辨认这一性质}

切换回平方函数，比较读数中的函数高度和切线斜率。平方函数只有个别位置会让二者相等；指数函数则处处相等。要检验“处处”这类命题，有限幅图只能帮助发现规律，最终仍需导数公式或微分方程的论证。



\paragraph{3. 有限时间的增长不能直接线性外推}

把步长调大，平均增长率不再等于起点的瞬时增长率。因为增长过程中存量本身持续增大，后面的增长也更快。复利中的细分极限把这种持续反馈纳入计算；因此一段时间的倍率是指数，而不是永远沿起点切线走。
